\section{Conclusions}\label{sec:conclusions}

We built a SimPy discrete-event model of the medium-voltage switchgear
warehouse at Schneider Electric's Manisa plant, drove it with four months of ZWM92
dispatch data, and compared six storage-assignment policies over 20
paired replications. The proposed Line-aware Slotting policy reduced
mean kitting lead time from 12.37 minutes per order under the current
SAP placement to 6.91 minutes --- a 44\% reduction --- and brought
reach-truck utilization from 21.7\% down to 1.3\%
(Table~\ref{tab:kpi_summary}). One-way ANOVA across the six policies is
strongly significant ($F = 47.1$, $p < 10^{-25}$), and the paired
$t$-test under common random numbers between Line-aware Slotting and
the current SAP baseline gives $p < 0.001$. The improvement did
not come from shorter walking paths --- average walk distance is about
92 metres per order under Line-aware Slotting versus 89 metres under the
SAP baseline. It came from removing reach-truck calls from the kit path,
which collapsed the operator queue that accounts for most of the lead
time.

\subsection*{Managerial implications}

For fast-moving materials --- those with meaningful pick frequency in the
ZWM92 log --- the model estimates that phasing in Line-aware Slotting
could cut kit preparation time by roughly 5.5 minutes per order. At
roughly 400 kit orders per working day, that translates to a substantial
daily saving in operator and reach-truck idle time, though the actual
gain in any given week depends on demand mix and fleet availability,
neither of which the model predicts. We recommend starting with a single
production line: identify the 50 to 100 materials that kit most
frequently to that line, move them to hand-reachable levels (the lower
three shelf levels, accessible without a truck) in the corridor that
serves the line, and leave the rest of the warehouse alone. The cost is
modest --- a one-time repallet of a small fraction of stock --- and if
the model reflects reality, the improvement should show up in lead-time
measurements within the first few weeks. The remaining five lines can
then be phased in as confidence grows.

\subsection*{Limitations}

Several modelling choices limit how directly the numbers should be read.
All element times come from a video time-motion study of the F400 kitting
line (2,319 micro-events) and were then extrapolated to the other eight
product families without per-line verification. Sensitivity analysis
shows that a $\pm$20\% change in operator pick time shifts lead time by
roughly the same proportion, so if F400 is not representative of a
slower or faster line the estimates will be off by that amount.

When a material appears in SAP at more than one bin, the model places it
at all decoded locations and the operator picks the nearest; 386
materials in the data have multiple bins, but the simulation does not
distinguish between forward and reserve assignments, so the interaction
between multi-bin splits and RT demand is only partially captured.

The simulation was also fitted and validated against the same ZWM92
dataset: the $\chi^2$ goodness-of-fit and the paired-$t$ volume checks
are internal-consistency tests, not holdout validation. The daily
dispatch rate gives some independent grounding --- the model generates
about 378 orders per simulated day against ZWM92's 396 active-day
average, a $-4.6$\% gap --- and the F400 video provides the timing
anchor, but no separate observation period was available.

Two further structural limits are worth stating plainly. Operator
congestion in shared aisles is not modelled: agents pass through each
other freely, so the model slightly under-estimates lead time under
high-load policies. And the scope is order picking only. Line-aware
Slotting concentrates fast movers at low levels, which will increase
replenishment frequency for those positions; reach-truck time spent on
replenishment is entirely outside the model. The reported RT utilization
figure of 1.3\% for Line-aware Slotting is therefore a picking-only
number, and total RT workload in a live deployment would be higher. For
completeness, only 750 of the active materials have fully decoded SAP
rack-bay-position codes; the remainder are either Kardex-routed or
placed by heuristic fallback, so the SAP baseline is best understood as
a SAP-plus-heuristic hybrid rather than a pure replay of the existing
layout.

\subsection*{Future work}

Four extensions would strengthen a follow-up study. Per-line video time
studies for SM6-36, Premset, MCSET, and the other families would replace
the F400-extrapolated constants with line-specific values and sharpen the
sensitivity bounds. A kit-aware co-location policy --- placing materials
that consistently appear in the same kit order near each other, using the
reconstructed BOMs from the ZWM92 grouping --- would complement
line-aware slotting and could reduce intra-order travel as well as
queueing. A physical pilot on one rack row, with lead-time measurements
before and after re-slotting, would provide holdout validation that the
simulation alone cannot supply. Finally, once more months of ZWM92 data
accumulate, a proper train/test split would allow a genuine out-of-sample
volume and frequency check, which is currently the study's main
evidential gap.

\subsection*{Sustainable development}

This project connects to three of the United Nations Sustainable
Development Goals \cite{sdgs}. SDG~9 (Industry, Innovation and
Infrastructure) is the most direct fit: replacing intuition-driven bin
assignment with a data-driven simulation analysis is evidence-based
industrial practice at essentially no capital cost --- no new equipment,
no expansion. SDG~12 (Responsible Consumption and Production) applies
because the proposed slotting makes better use of the storage positions
and handling equipment that already exist; less reach-truck idle time
means less energy consumed per kit order, and avoiding unnecessary
building expansion preserves materials and land. SDG~8 (Decent Work and
Economic Growth) matters on the operator side: fewer high-level picks
means fewer minutes standing idle waiting for a truck, and the
lower-level hand picks that replace them are ergonomically less demanding
than mast-elevated retrieval. None of these are transformative claims.
They are modest, honest improvements that follow directly from what the
model found.
